
%! Author = chongsan
%! Date = 2026/7/9
% 仅优化排版格式，正文所有文字、公式、内容完全未做任何修改

\documentclass[11pt,a4paper]{article}
\usepackage[margin=1in]{geometry}
\usepackage{amsmath,amssymb,amsthm} % amsthm提前加载

% 【核心修复】定理环境放导言区，英文内部标识，中文显示
\newtheorem{defi}{定义}[section]
\newtheorem{thm}{定理}[section]
\newtheorem{lem}{引理}[section]
\newtheorem{cor}{推论}[section]

\usepackage{tabularx}
\usepackage{enumitem}
\usepackage[UTF8]{ctex}
\usepackage{unicode-math}
\usepackage{indentfirst}
\usepackage{caption}
\usepackage{hyperref} % hyperref必须所有宏最后

% 全局列表统一格式
\setlist[enumerate]{leftmargin=*,itemsep=0.3em}
\setlist[itemize]{leftmargin=*,itemsep=0.3em}

% 作者标注定义
\newcommand{\authortarjan}{Robert E. Tarjan \textsuperscript{\dag}}
\newcommand{\authorzwick}{Uri Zwick \textsuperscript{\ddag}}

\begin{document}

% 阅读笔记区块（原文一字未改，仅排版美化）
\noindent 2026年7月9日
\begin{quote}
已通读豆包AI翻译的文档，仍有一些不解之处，例如：
\begin{enumerate}
    \item 第6部分的常驻空间上界 $N+(2r-1)+2(r-1)B+B=N+O(rN^{1/r})$ 的这个系数2是什么
    \item 下标访问原生$O(r)$最坏时间是为什么，按我的理解顶多是二级指针，$O(1)$复杂度
    \item 第6.2部分信用额度是什么鬼第一次听说，再看看
    \item 第8部分看得有点云里雾里的，怀疑AI有删减
\end{enumerate}
此外，还有一种局部看得懂，整合起来又有点晕晕的奇妙感觉

大致搞懂了它的数据结构构造（大成功？），复杂度分析也大致明白了

趣事：第6.1部分推论6.2中我还在想$r=\log N$时常驻空间为什么不是$O(\log N\cdot N^{1/\log N})$，突然意识到$N^{1/\log N}$是个常数
\end{quote}
以下为AI翻译的文档：
\vspace{1em}

% 论文标题、作者、摘要
\title{最优可变长数组}
\author{\authortarjan \quad \authorzwick}
\maketitle

\begin{abstract}
可变长数组支持在数组尾部（或两端）增删元素，同时能够根据下标常数时间访问任意元素。由于数组元素数量会动态变化，要实现空间高效的可变长数组，必须配套动态内存管理机制。经典\textbf{倍增扩容}方案可实现规模为$N$的数组仅占用$O(N)$空间，且单次操作均摊时间$O(1)$，甚至最坏时间$O(1)$。Sitarski、Brodnik等人提出了更优方案：存储规模$N$的数组仅需$N+O(\sqrt{N})$空间，单次操作仍保持$O(1)$时间；Brodnik等人还给出证明，该空间复杂度是理论下界。

本文区分两类空间开销：一是数组常态存储、元素访问所需的常驻空间；二是扩容/缩容操作过程中临时占用的额外空间。对任意整数$r\ge2$，我们证明：若扩容、缩容时可短暂使用$N+O(N^{1-1/r})$临时空间，则常态存储规模$N$的数组仅需$N+O(N^{1/r})$常驻空间；下标访问最坏时间$O(1)$，增删操作均摊时间$O(r)$。

我们通过一套抽象\textbf{扩容博弈}模型做精确分析，证明：若某类通用数据结构仅允许$N+O(N^{1/r})$常驻存储空间，即便只支持尾部添加与下标查询，尾部新增操作的均摊代价下界也为$\Omega(r)$。除非$r=2$，否则增删操作无法做到最坏时间$O(1)$。
\end{abstract}

\section{引言}
数组与可变长数组是应用最广泛的数据结构。令人意外的是，本文将证明，关于它们仍存在大量未被充分挖掘的理论结论。我们提出一套简洁的全新可变长数组实现方案，多数场景下内存利用率优于此前所有设计。本文视角偏向理论，但给出的实现方案也具备实际工程价值。

规模为$N$的数组是存储序列$a_0,a_1,\dots,a_{N-1}$的抽象数据结构：对任意$0\le i<N$，下标$i$对应元素可在常数时间读取/修改。数组最基础的实现方式是连续内存块，每块内存存储一个元素或元素指针，依托机器随机访问能力实现常数时间读写；但原生数组的容量固定。

可变长数组（又称动态数组、动态表）允许数组容量动态增减，同时保证任意下标元素可常数时间读写。单纯缩小数组很简单：只需不再使用已分配内存的后半段，但这种方式内存利用率极低；扩容则更棘手——数组连续内存块末尾的内存大概率已被其他数据占用，因此可变长数组必须依赖动态内存管理。

为简化讨论，本文大部分内容仅讨论\textbf{单端可变长数组}（仅在数组尾部增删元素）：数组规模从$N$增至$N+1$时，在末尾新增元素$a_N$；规模从$N$缩减至$N-1$时，丢弃末尾元素$a_{N-1}$。这类可变长数组是栈的泛化结构。本文所有结论均可拓展至两端均可增删的双端可变长数组（即双端队列deque）：只需用两个单端可变长数组拼接实现即可。

教材中经典的可变长数组实现是\textbf{倍增扩缩容}（见Cormen等人《算法导论》17.4节）：初始分配一块定长内存；内存填满时，分配一块两倍大小的新内存，将全部元素拷贝过去并释放旧内存；当内存使用率低于1/4时，分配一半大小的新内存拷贝元素、释放旧内存。均摊分析可证明，单次扩容/缩容操作均摊代价$O(1)$，元素读写最坏时间$O(1)$；借助后台渐进重构技术，甚至能将扩缩容优化为最坏时间$O(1)$。

维基百科记载，Java \texttt{ArrayList}、Python \texttt{PyListObject}、C++ STL \texttt{vector}等主流容器均采用几何扩缩容策略，扩容系数取1.25~2。

倍增方案仅保证$O(N)$空间复杂度，虽能满足多数场景，但纯扩容场景下最多会浪费50\%已分配内存。更一般地，若扩容系数为$1+\alpha$，内存闲置比例可达$\frac{\alpha}{\alpha+1}$；若缩小$\alpha$降低闲置占比，单次扩容均摊代价会上升至$O(\frac{1+\alpha}{\alpha})$。

一个核心问题：能否构造仅占用$N+o(N)$空间、且所有操作$O(1)$时间的可变长数组？Sitarski与Brodnik等人给出肯定答案：设计实现仅需$N+O(\sqrt{N})$空间，我们将在第3节详细介绍。Brodnik等人同时证明该方案存在下界——任意可变长数组实现，在某些时刻必然占用$N+\Omega(\sqrt{N})$空间，第4节会复述其简洁证明。

本文核心创新：区分\textbf{常驻存储空间}（常态存放数组、支持下标读写所需空间）与\textbf{扩缩容临时空间}（扩容/缩容过程中短暂占用的额外内存）。很多场景下，我们愿意牺牲少量临时内存，换取常态存储更高的紧凑度；例如程序中存在大量可变长数组，但同一时刻仅一个数组执行扩缩容。

我们首个突破Brodnik等人下界的实现方案：常态存储仅$N+O(N^{1/3})$空间，仅在扩缩容时临时占用$N+O(N^{2/3})$空间；元素读写最坏时间$O(1)$，增删均摊$O(1)$。同时我们证明该方案渐近最优：任何仅用$N+O(N^{1/3})$常驻空间的实现，扩缩容时必然需要$N+\Omega(N^{2/3})$临时空间。

更一般地，对任意$r\ge2$：
\begin{enumerate}
    \item 常态存储开销：$N+O(N^{1/r})$；
    \item 扩缩容临时开销：$N+O(N^{1-1/r})$；
    \item 下标访问最坏时间$O(1)$；
    \item 尾部增删操作均摊时间$O(r)$。
\end{enumerate}
由于无法使用后台重构技术，除非$r=2$，否则$O(r)$的均摊代价无法优化为最坏时间$O(1)$。$r=2$对应Sitarski、Brodnik的经典方案。

\subsection{论文章节结构}
\begin{enumerate}
    \item 第2节：精确形式化问题定义与计算模型；
    \item 第3节：梳理前人相关工作，重点介绍Sitarski的HAT结构、Brodnik等人的方案；
    \item 第4节：复述Brodnik的下界证明，并给出我们的拓展下界；
    \item 第5节：给出$r=3$的完整构造（常驻$O(N^{1/3})$、临时$O(N^{2/3})$）；
    \item 第6节：对任意$r\ge2$给出通用无递归实现，证明增删均摊$O(r)$；
    \item 第7节：提出一套结构变换技巧，将常驻空间从$O(rN^{1/r})$优化至$O(N^{1/r})$，同时不破坏常数访问、$O(r)$均摊增删；
    \item 第8节：抽象出\textbf{扩容博弈}模型并完整求解，用于证明下界；
    \item 第9节：依托扩容博弈，证明仅使用$N+O(N^{1/r})$常驻空间时，尾部新增操作均摊代价下界$\Omega(r)$；
    \item 第10节：总结与展望。
\end{enumerate}

\section{可变长数组形式化定义}
可变长数组作为抽象数据类型，支持以下操作：
\begin{enumerate}
    \item $\text{Array}()$：创建空数组并返回；
    \item $\text{A.Length}()$：返回数组当前元素个数；
    \item $\text{A.Get}(i)$：返回下标$i$对应的元素$a_i$，要求$0\le i<\text{A.Length}()$；
    \item $\text{A.Set}(i,a)$：将下标$i$元素修改为$a$，要求$0\le i<\text{A.Length}()$；
    \item $\text{A.Grow}(a)$：数组长度$+1$，末尾新增元素$a$；
    \item $\text{A.Shrink}()$：数组长度$-1$，丢弃末尾元素。
\end{enumerate}

本文仅允许在数组尾部增删元素。若支持任意位置插入删除（并同步平移后续元素下标），问题难度会大幅提升；此时所有操作最优时间界为$O(\frac{\log N}{\log\log N})$，且该界紧（Dietz、Fredman\&Saks）。另有一类方案：下标访问最坏$O(r)$，插入删除均摊$O(N^{1/r})$（Goodrich、Joannou等人）。再次强调：本文仅讨论尾部增删场景。

我们假设内存管理器支持分配、释放任意长度的定长数组（类比C/C++的\texttt{malloc}/\texttt{free}），简化起见，单次分配/释放操作代价为常数；即便分配长度$N$的内存块代价为$O(N)$，本文所有均摊复杂度结论依然成立。本文不讨论内存管理器底层实现。

可变长数组实现由若干动态分配的定长内存块（简称块）组成，分为三类：
\begin{enumerate}
    \item 数据块：存储数组元素，每个元素占一个机器字；
    \item 索引块：存储指向其他块的指针，每个指针占一个机器字；
    \item 辅助块：存放数据结构元信息，如各块长度。
\end{enumerate}

本文所有实现均遵循该框架。若数组当前规模为$N$，则全部数据块总容量至少为$N$（第4节定理4.1证明后补充讨论）。除主索引块外，所有块必须被某索引块内的指针引用，否则无法访问。数据结构总空间为所有已分配块的长度之和。

\begin{defi}[$(s(N),t(N))$实现]
设$s(N),t(N)$为单调不减函数。若某可变长数组满足：
\begin{enumerate}
    \item 存储规模$N$的数组时，常驻总空间不超过$N+s(N)$；
    \item 执行扩容/缩容操作过程中，瞬时总空间不超过$N+t(N)$；
\end{enumerate}
则称其为$(s(N),t(N))$实现。
\end{defi}

\section{前人相关工作}
设计空间高效可变长数组属于紧凑数据结构领域，Raman、Munro等人已有大量综述，但仅有少数文献专门研究本文定义的单端可变长数组。

\subsection{基础单块实现}
基础实现仅用一块连续定长内存承载数组：内存填满时，分配更大内存、拷贝全部元素、释放旧内存；内存空置过多时，分配更小内存拷贝元素。实现可自由选择新内存块的尺寸。
\begin{enumerate}
    \item 极致省常驻空间方案：每次增删都重新分配恰好等于元素数量的内存。常驻空间$N+O(1)$（仅存长度与内存指针），但单次增删均摊代价$\Omega(N)$；扩缩容时需同时保存新旧两块内存，临时空间$2N+O(1)$，仅适合极少扩缩容的场景。
    \item 几何扩缩容（工程主流）：固定参数$\alpha>0$，内存填满时分配$(1+\alpha)N$大小新块；内存元素仅占当前块容量$\frac{1}{(1+\alpha)^2}$时，分配$\frac{N}{1+\alpha}$大小新块拷贝元素。简单计算可得：扩容均摊代价$\frac{1+\alpha}{\alpha}$，缩容均摊代价$\frac{1}{\alpha}$。合理选取$\alpha$可平衡空间与时间，适配绝大多数工程场景。
\end{enumerate}

\subsection{Sitarski的哈希数组（HAT）}
Sitarski提出名为HAT（哈希数组树）的结构，但命名存在误导——不含哈希，本质分层链表。

HAT维护长度为$B$的索引块，满足$\sqrt{N}\le B<4\sqrt{N}$；配套$\lceil N/B\rceil$或$\lceil N/B\rceil+1$个尺寸$B$的数据块。$N$个元素顺序存入前$\lceil N/B\rceil$块；$N$无法整除$B$时最后一块半满；存在$\lceil N/B\rceil+1$块则为空。索引块第$i$项存储第$i$个数据块指针。

总常驻空间：$N+3B+2=N+O(\sqrt{N})$。$3B$为索引块+两块闲置数据块开销。下标$i$元素位于第$\lfloor i/B\rfloor$块、偏移$i\bmod B$，读写最坏$O(1)$；$B$取2的幂时可通过位运算快速计算下标。

扩容逻辑：$\sqrt{N}<B$即$N<B^2$时，末尾块未满直接写入；已满但存在空块则写入空块；无空块则分配新$B$块。分配代价常数则扩容最坏$O(1)$；分配代价$O(B)$时扩容均摊$O(1)$。

满载$N=B^2$时将$B$翻倍重构；初始$B$为2的幂则重构后不变。重构$O(N)$拷贝代价分摊至$\Omega(N)$次扩容，均摊$O(1)$。

缩容对称：仅当数组元素总量降至$B^2/16$（$B=4\sqrt{N}$）、末尾两块全空时才折半重构，避免频繁扩缩震荡。两次重构间至少$\Omega(N)$次操作，缩容均摊$O(1)$。

重构仅需$O(\sqrt{N})$临时空间，新旧块逐块分配释放，无需完整保存两份数组。

可通过维护两种尺寸数据块+后台渐进重构，转化为最坏时间界版本，但结构复杂度大幅提升；下一节方案无需全局重构、元素全程不移动。

\subsection{Brodnik等人的分层块结构}
Brodnik独立提出分层方案，常驻$N+O(\sqrt{N})$，扩缩容最坏$O(1)$（分配常数时间）；优势：无全局重构、元素不移动。两种变体：
\begin{enumerate}
    \item 基础变体：数据块尺寸依次$1,2,3,\dots$；前$k$块总容量$\approx k^2/2$，$k=\lceil\frac{\sqrt{8N+1}-1}{2}\rceil$；末块半满，预留一块空块。索引块尺寸$\Theta(\sqrt{N})$，自身朴素动态数组实现。有空位直接写入；无空位分配新块，索引块满则扩容拷贝指针，后台拷贝实现最坏$O(1)$。下标块编号$\ell=\lceil\frac{\sqrt{8i+1}-1}{2}\rceil$，块内偏移$i-\frac{\ell(\ell-1)}{2}$，读写$O(1)$但需要开平方运算。
    \item 无平方根变体：所有数据块尺寸为2的幂，引入多层超块结构，利用移位运算替代开平方定位下标，规避浮点与开方开销，整体空间复杂度依旧$N+O(\sqrt{N})$，索引总开销仍为$\Theta(\sqrt{N})$级别。
\end{enumerate}

\section{Brodnik下界定理及其拓展}
本节复述原始下界与分离空间拓展下界。
\begin{thm}
任意仅支持扩容、下标访问的可变长数组，在某些时刻必然占用$N+\Omega(\sqrt{N})$空间，$N$为当前元素总数。
\end{thm}
\begin{proof}
完成$N$次扩容后设使用$k$块连续内存；全部数据块容量至少$N$，每块需一条指针，指针总开销$k$。
\begin{itemize}
    \item 若$k\ge\sqrt{N}$：总空间$\ge N+\sqrt{N}$，下界成立；
    \item 若$k<\sqrt{N}$：必存在尺寸$\ell>\sqrt{N}$的块。该块在数组规模$N'\le N$时分配，分配瞬间总空间$N'+\ell>N'+\sqrt{N'}$，下界成立。
\end{itemize}
证毕。
\end{proof}

下界两条基础假设：
\begin{enumerate}
    \item $N$个独立$w$比特元素至少占用$N$机器字；
    \item $k$条指针占用$k$机器字（地址空间$2^w$）。
\end{enumerate}
地址空间$2^{cw}(0<c<1)$时指针开销为$ck$，仅常数变化，不改变渐近阶。

\begin{thm}[$(s(N),t(N))$分离下界]
任意仅支持扩容、访问的$(s(N),t(N))$实现，满足$s(N)\cdot t(N)\ge N$。
\end{thm}
\begin{proof}
$N$个元素分散至多$s(N)$块，鸽巢原理存在块尺寸$\ge N/s(N)$。该块分配时刻瞬时额外空间$\ge N/s(N)$；$t(N)$单调不减，$t(N)\ge N/s(N)$，变形得$s(N)t(N)\ge N$。
证毕。
\end{proof}
\begin{cor}
对任意整数$r\ge1$，常驻仅$N+O(N^{1/r})$的实现，扩容瞬时必有$N+\Omega(N^{1-1/r})$额外空间。
\end{cor}
后文证明该下界是紧的（存在实现刚好匹配该空间权衡）。

\section{\texorpdfstring{$r=3$}{r=3}的极简实现：\texorpdfstring{$(O(N^{1/3}),O(N^{2/3}))$}{(O(N^1/3),O(N^2/3))}方案}
构造常驻$O(N^{1/3})$、临时$O(N^{2/3})$，访问最坏$O(1)$，增删均摊$O(1)$。

核心改造HAT结构：外层大块尺寸$N^{2/3}$、一级索引开销$N^{1/3}$；半满尾部块嵌套一层HAT，嵌套层额外开销仅$O((N^{2/3})^{1/2})=O(N^{1/3})$。合并拷贝瞬时需要$\Omega(N^{2/3})$临时内存，匹配推论下界。也可基于Brodnik分层思路构造等价方案。

无递归完整描述：维护参数$N^{1/3}\le B\le4N^{1/3}$，分两类内存块：
\begin{enumerate}
    \item 大块：尺寸$B^2$，总数量约$N/B^2$，所有大块永久填满；
    \item 小块：尺寸$B$，最多允许$2B$个，至多一块部分填充、至多一块完全闲置空块。
\end{enumerate}
配套两级索引分别管理大块、小块；当累计$B$个完整小块时可合并为一块大块；无可用小块时拆分一块大块生成小块。所有访问最坏$O(1)$，增删均摊$O(1)$。

\section{对任意\texorpdfstring{$r\ge2$}{r≥2}的通用\texorpdfstring{$(O(rN^{1/r}),O(N^{1-1/r}))$}{(O(rN^1/r),O(N^1-1/r))}实现}
泛化$r=3$构造，无递归分层设计，常驻额外$O(rN^{1/r})$，扩容临时峰值$O(N^{1-1/r})$。

设定参数$N^{1/r}\le B<4N^{1/r}$，所有数据块尺寸为$B$的幂$B,B^2,\dots,B^{r-1}$。仅扩容场景每层最多$B$块；同时支持增删时放宽至每层最多$2B$块（冗余基数计数器，避免频繁震荡）。
\begin{itemize}
    \item $B^2,\dots,B^{r-1}$尺寸块永久填满；
    \item 仅最末尾$B$尺寸块允许部分填充。
\end{itemize}

扩容逻辑：末尾$B$块有空位直接写入；已满且总数不足$2B$则分配新$B$块；若达到$2B$个完整小块，则合并$B$个小块生成高一阶大块，逻辑等价基数$B$进位计数器。

当总元素$N=B^r$时，将参数$B$翻倍全局重构；配套$r-1$层独立索引块。常驻空间上界$N+(2r-1)+2(r-1)B+B=N+O(rN^{1/r})$；合并$B^{r-1}$尺寸大块时临时空间峰值$O(N^{1-1/r})$。

原生未优化时下标访问最坏$O(r)$；6.3节位运算优化后可达$O(1)$最坏时间。增删均摊代价$O(r)$：每个元素先存入$B$小块，逐层合并提升块阶，全局重构分摊后仅增加常数开销。

缩容适配：每层上限$2B$块；达到$2B$完整小块时合并；无可用小块则拆分最小阶大块，均摊代价依旧$O(r)$。

\subsection{6.1 实现伪代码说明}
$n_i$代表尺寸$B^i$的数据块数量，$A[i]$为对应索引块；$n_0$为最后一块$B$块内已存元素数量。
\subsubsection{Grow(a) 扩容流程}
\begin{enumerate}
    \item 若$N=B^r$，执行$\text{Rebuild}(2B)$，$B$翻倍重构；
    \item 若$n_1=2B,n_0=B$（所有$B$块填满），调用$\text{Combine-Blocks}$合并低层块；
    \item 无空余位置则分配全新$B$块，重置$n_0=0$；
    \item 将新元素写入末尾$B$块空位，$n_0,N$自增1。
\end{enumerate}
$\text{Combine-Blocks}$：找到最小未满层$k$，自底向上每层合并$B$个小块为一块高阶大块，调整索引指针、释放旧内存块。

\subsubsection{Shrink() 缩容流程}
\begin{enumerate}
    \item 若$N=(B/4)^r$，执行$\text{Rebuild}(B/2)$，$B$折半重构；
    \item 若无$B$尺寸块，调用$\text{Split-Blocks}$拆分大块；
    \item 删除末尾元素，$n_0,N$自减1；若$n_0=0$，释放空$B$块。
\end{enumerate}
$\text{Split-Blocks}$：找到最小非空大块，拆分为多层小块，分配新内存、拷贝元素、释放原大块。

辅助函数：
\begin{itemize}
    \item $\text{Allocate}(\ell)$：分配长度$\ell$内存块，返回指针；
    \item $\text{Deallocate}(X)$：释放内存块$X$；
    \item $\text{Copy}(X,x,Y,y,\ell)$：从$X$偏移$x$拷贝$\ell$个元素至$Y$偏移$y$；
    \item $\text{Rebuild}(B')$：以新参数$B'$全局重构整套分层结构。
\end{itemize}

\begin{thm}
该实现为$(O(rN^{1/r}),O(N^{1-1/r}))$可变长数组；扩容、缩容均摊$O(r)$，下标读写最坏$O(1)$。
\end{thm}

\begin{cor}
取$r=\log N$，存在实现常驻额外$O(\log N)$，增删均摊$O(\log N)$，访问最坏$O(1)$。
\end{cor}

\subsection{6.2 增删操作均摊代价证明（引理6.3）}
定义单次操作代价为元素拷贝赋值次数。
\begin{enumerate}
    \item 扩容：基础写入代价1；每个底层元素分配$r-1$单位信用额度，每向上合并一层消耗1信用，合并操作均摊代价归零，单次扩容基础均摊$O(r)$；
    \item 缩容：每次缩容为所有层级增加固定信用；拆分大块的拷贝开销由长期累积信用覆盖，拆分均摊代价归零，单次缩摊$O(r)$；
    \item 全局重构：两次重构间至少执行$\Omega(N)$次增删，重构总拷贝$N$可完全分摊，不提升渐近阶。
\end{enumerate}
综上扩容、缩容均摊代价$O(r)$。

\subsection{6.3 \texorpdfstring{$O(1)$}{O(1)}下标访问优化（引理6.4）}
假设CPU内置最高置位$\text{msb}$常数查询指令，对结构小幅改造：
\begin{enumerate}
    \item 令$B=2^b$（2的幂次）；
    \item 总元素$N$可表示冗余$B$进制：$N=\sum_{j=0}^{r-1}n_j B^j,\;0\le n_j\le2B$；
    \item 每层块计数压缩存入两个机器字$N^0,N^1$，移位与按位运算快速提取；
    \item 预计算每层容量前缀和$N_k$，代表尺寸不超过$B^{k-1}$的总元素数量。
\end{enumerate}

访问下标$i$步骤：
\begin{enumerate}
    \item 计算相对末尾偏移$x=N-1-i$；
    \item 调用$\text{msb}$定位$x$最高进制位，确定目标块阶$k$；
    \item 移位、按位与替代取模/除法，快速计算块内偏移，全程常数操作。
\end{enumerate}
若无原生$\text{msb}$指令，现有算法耗时$O(\log\log w)$，不改变渐近$O(1)$；工程场景$r$通常极小，直接逐层遍历也具备实用性。

\section{数据结构变换优化常驻空间}
当$r(N)$随$N$缓慢增长时，可将常驻额外空间从$O(rN^{1/r})$降至$O(N^{1/r})$，同时保留$O(r)$均摊、$O(1)$访问性能。
\begin{defi}[标准$(s(N),t(N))$实现]
满足两条约束：
\begin{enumerate}
    \item 分配新块时按顺序拷贝若干旧块内容，拷贝完成立刻释放旧块；
    \item 拷贝操作结束后，常驻空间回归$N+s(N)$级别。
\end{enumerate}
本文全部实现均属于标准实现。
\end{defi}

\begin{lem}
任意标准$(s(N),O(N))$实现，可变换得到$(s(N)+O(N/t(N)),s(N)+O(t(N)))$新实现，访问、均摊复杂度不变。
\end{lem}
思路：超大块拷贝拆分多段，每段不超过$t(N)$；新增指针开销$O(N/t(N))$，瞬时空间峰值变为$s(N)+O(t(N))$。

\begin{thm}
对满足$r(N)\le\frac{1}{2}\frac{\log N}{\log\log N}$的$r$，存在$(O(N^{1/r}),O(N^{1-1/r}))$实现：访问最坏$O(1)$，增删均摊$O(r)$。
\end{thm}
证明思路：对$r'=2r$的基础构造应用变换，取$t(N)=N^{1-1/r}$；约束下$rN^{1/(2r)}\le N^{1/r}$，常驻开销简化为$O(N^{1/r})$。

\section{扩容博弈：下界抽象模型}
为证明常驻空间受限场景下扩容均摊下界，抽象单人扩容博弈模型，刻画所有标准实现的合并行为。

\subsection{8.1 \texorpdfstring{$(N,k,\ell)$}{(N,k,l)}扩容博弈定义}
参数：总插入$N$、最大子数组数量$k$、单块最大空位$\ell$。
初始$k$个空数组$A_1\dots A_k$；元素存储顺序倒置；$a_i$为$A_i$内元素个数。约束：
\begin{enumerate}
    \item 空数组全部集中在最左侧；
    \item 仅首个非空数组允许保留至多$\ell$空位，其余全部填满；
    \item 总闲置开销等价$k+\ell$，对应数据结构常驻额外空间。
\end{enumerate}

单次插入三类操作：
\begin{enumerate}
    \item 首块存在空位直接写入，代价1；
    \item 无空位但存在空数组，分配长度$\ell+1$新块存入元素，代价1；
    \item 所有子数组填满：选定$i$，合并前$i$块为新$A_i$，代价$1+\sum_{j=1}^i a_j$。
\end{enumerate}
记$C_{N,k,\ell}$为$N$次插入最小总代价；单次均摊$A_{N,k,\ell}=C_{N,k,\ell}/N$。

\subsection{8.2 \texorpdfstring{$\ell=0$}{l=0}归约至无空位场景}
\begin{lem}
若$N$可被$\ell+1$整除，则$C_{N,k,\ell}=(\ell+1)C_{N,k,0}$，单次均摊代价相等。
\end{lem}
直观理解：每填满$\ell+1$个空位等价完成一轮块填充，可把每$\ell+1$次插入打包为$\ell=0$博弈单元。

\subsection{8.3 \texorpdfstring{$\ell=0$}{l=0}博弈精确求解}
记$C_{N,k}=C_{N,k,0}$；状态集合$P_{N,k}$满足总元素$\sum a_i=N$、空数组前置。
\begin{thm}
\begin{enumerate}
    \item 若$\binom{n+k-1}{k}-1\le N\le\binom{n+k}{k}-1$，总最小代价$C_{N,k}=(N+1)n-\binom{n+k}{k+1}$；
    \item 区间内每次新增元素增量代价恒为$n$：$C_{N,k}-C_{N-1,k}=n$；
    \item 最优状态充要条件：$\binom{n+i-2}{i}\le a_i\le\binom{n+i-1}{i},\forall i\in[k]$。
\end{enumerate}
\end{thm}
\begin{cor}
\begin{enumerate}
    \item 当$N=\binom{n+k}{k}-1$时，总代价$C_{N,k}=\frac{kn}{k+1}(N+1)$；
    \item 区间内均摊代价下界$A_{N,k}\ge\frac{1}{2}(n-1)$。
\end{enumerate}
\end{cor}

\subsection{8.4 最优策略：二项式计数器}
唯一最优插入序列对应二项式计数器；维护计数$b_1\dots b_k$，每次寻找最小$i$满足$b_i<b_{i+1}$并合并进位；全部$b_i=n$时策略唯一。

\subsection{8.5 博弈规则拓展}
允许任意区间无序合并不会降低最小总代价；拓展博弈完整覆盖本文所有块合并、拆分逻辑。

\section{扩容操作均摊代价下界}
\begin{thm}
任意标准可变长实现，常驻仅$O(rN^{1/r})$、$r=O(\log N)$时，扩容均摊代价下界$\Omega(r)$。
\end{thm}
证明思路：取博弈参数$k=\Theta(rN^{1/r})$，利用二项式系数放缩证明博弈最优均摊至少线性于$r$；所有标准实现行为等价博弈合法操作，故数据结构下界为$\Omega(r)$。

\section{总结}
本文提出一族分层可变长数组实现，完整刻画**常驻空间、扩容临时空间、增删均摊时间**三者最优权衡关系，同时保证$O(1)$下标访问，从理论层面几乎完整解决该基础数据结构的时空权衡问题。

现有下界仅适用于本文定义的标准实现；引入元素、指针不可压缩假设后，下界可拓展至全部实现，无本质证明障碍。

扩容博弈具备独立理论价值；未来可探索无需完整闭式解的归纳式下界证明。

% 参考文献
\begin{thebibliography}{17}
\bibitem[1]{bille2017fast} Philip Bille, Anders Roy Christiansen, and Inge Li Gørtz. Fast dynamic arrays. In Proc. of 25th ESA, volume 87 of LIPIcs, pages 16:1–13, 2017.
\bibitem[2]{brodnik1999resizable} Andrej Brodnik, Svante Carlsson, Erik D. Demaine, J. Ian Munro, and Robert Sedgewick. Resizable arrays in optimal time and space. In Proc. of 6th WADS, pages 37–48, 1999.
\bibitem[3]{brodnik1997trans} Andrej Brodnik, Peter Bro Miltersen, and J. Ian Munro. Trans-dichotomous algorithms without multiplication. In Proc. of 4th WADS, volume 1272 of Lecture Notes in Computer Science, pages 426–439. Springer, 1997.
\bibitem[4]{cormen2009intro} Thomas H. Cormen, Charles E. Leiserson, Ronald L. Rivest, and Clifford Stein. Introduction to algorithms. MIT press, 3rd edition, 2009.
\bibitem[5]{dietz1989optimal} Paul F. Dietz. Optimal algorithms for list indexing and subset rank. In Proc. of 1st WADS, volume 382 of Lecture Notes in Computer Science, pages 39–46. Springer, 1989.
\bibitem[6]{fredman1989cell} Michael L. Fredman and Michael E. Saks. The cell probe complexity of dynamic data structures. In Proc. of 21st STOC, pages 345–354. ACM, 1989.
\bibitem[7]{fredman1993surpass} Michael L. Fredman and Dan E. Willard. Surpassing the information theoretic bound with fusion trees. J. Comput. Syst. Sci., 47(3):424–436, 1993.
\bibitem[8]{goodrich1999tiered} Michael T. Goodrich and John G. Kloss II. Tiered vectors: Efficient dynamic arrays for rankbased sequences. In Proc. of 6th WADS, pages 205–216, 1999.
\bibitem[9]{joannou2011empirical} Stelios Joannou and Rajeev Raman. An empirical evaluation of extendible arrays. In Proc. of 10th SEA, volume 6630 of Lecture Notes in Computer Science, pages 447–458. Springer, 2011.
\bibitem[10]{kaplan2022simulate} Haim Kaplan, Robert E. Tarjan, Or Zamir, and Uri Zwick. Simulating a stack using queues. In Proc. of 33rd SODA, pages 1901–1924, 2022.
\bibitem[11]{katajainen2016worst} Jyrki Katajainen. Worst-case-efficient dynamic arrays in practice. In Proc. of 15th SEA, volume 9685 of Lecture Notes in Computer Science, pages 167–183. Springer, 2016.
\bibitem[12]{knuth2009taocp} Donald E Knuth. The Art of Computer Programming, Volume 4, Fascicle 1: Bitwise Tricks \& Techniques; Binary Decision Diagrams. Addison-Wesley, 2009.
\bibitem[13]{munro2018succinct} J. Ian Munro and S. Srinivasa Rao. Succinct representation of data structures. In Handbook of Data Structures and Applications, pages 595–610. Chapman and Hall/CRC, 2018.
\bibitem[14]{raman2001succinct} Rajeev Raman, Venkatesh Raman, and S. Srinivasa Rao. Succinct dynamic data structures. In Proc. of 7th WADS, volume 2125 of Lecture Notes in Computer Science, pages 426–437. Springer, 2001.
\bibitem[15]{sitarski1996hat} Edward Sitarski. Algorithm alley: HATs: Hashed array trees. Dr. Dobb’s Journal of Software Tools, 21(9):107–11, September 1996.
\bibitem[16]{wiki-dynarray} Wikipedia contributors. Dynamic array - Wikipedia, the free encyclopedia, 2022.
\bibitem[17]{wiki-hat} Wikipedia contributors. Hashed array tree - Wikipedia, the free encyclopedia.
\end{thebibliography}

\vspace{2em}
\section{关键术语对照表}
\begin{tabularx}{\linewidth}{|l|X|}
    \hline
    英文原文 & 标准中文译法 \\
    \hline
    resizable array & 可变长数组 \\
    dynamic array / vector & 动态数组 \\
    amortized time & 均摊时间 \\
    worst-case time & 最坏时间 \\
    block & 内存块 \\
    data block & 数据块 \\
    index block & 索引块 \\
    doubling technique & 倍增扩容法 \\
    geometric expansion & 几何扩缩容 \\
    HAT (Hashed Array Tree) & 哈希数组树 \\
    succinct data structure & 紧凑数据结构 \\
    lower bound / upper bound & 下界 / 上界 \\
    cell probe model & 单元探测模型 \\
    binomial counter & 二项式计数器 \\
    growth game & 扩容博弈 \\
    temporary space & 临时空间 \\
    resident space & 常驻空间 \\
    rebuild & 全局重构 \\
    combine / split blocks & 块合并 / 块拆分 \\
    \hline
\end{tabularx}
以上为AI翻译的文档。

\vspace{1em}
\noindent 2026年7月9日
\begin{quote}
    第二次通读豆包AI翻译的文档（你问我为什么不去读原文？因为要做AI review，需要重新读AI翻译文档）

    第二次通读仍有一些问题，例如：
    \begin{enumerate}
        \item 第1部分引言中若扩容系数为$1+\alpha $，内存闲置比例可达$\frac{\alpha}{\alpha+1}$而不是$\frac{1}{\alpha+1}$（不知是AI译错还是写错）
        \item 第3.2部分总常驻空间依旧惊现神秘常数系数2，推测为索引块指针占用（但为什么是2呢？）
        \item 第3.3部分无平方根变体没有细讲，不知道复杂度怎么做到$O(N+\sqrt N)$ (我寻思新块的大小不是比之前所有块加起来还大吗)
        \item 第5部分为什么小块至多为2B个呢？我寻思小块有B个的时候不就可以合并成$B^2$的大块了吗
        \item 第7部分引理7.1有什么用吗？倒不如说一整个第7部分都很怪
    \end{enumerate}
    发现一些有意思的相关研究：
    \begin{enumerate}
        \item 第2部分中提到的支持任意位置插入删除数组实现（不是链表！？）
        \item 第3.2部分双尺寸块+后台重构转为最坏时间版本
    \end{enumerate}
    一些小小的思考：
    \begin{enumerate}
        \item 第3.2部分为什么$N=\frac{B^2}{16}$才重构缩容而不是$N=\frac{B^2}{4}$?\\
        推测为为了保证缩容后增容的空间，若$N=\frac{B^2}{4}$就缩容，那缩容后$N=\frac{B^2}$,符合增容条件再次增容，产生振荡
        \item 第3.2部分重构的缩容增容的系数都是B乘2 B除2，为什么依旧是神秘系数2?\\
        推测是为了防止过度重构吧，不过我觉得没什么所谓就是了，顶多就是改变一下复杂度前的常数的事。不过2也是一个很特殊的数字呢，2进制，2分法，2叉树，最小二乘法。。。2跟算法也有密不可分的联系呢。
    \end{enumerate}
    对第一次阅读问题的解决：
    \begin{enumerate}
        \item N+(2r−1)+2(r−1)B+B的(2r−1)系数2好像真是指针大小，2(r−1)B系数2好像是最多支持每层有2B块\\
        具体而言该公式的组成如下：\\
        N:有效数据占用\\
        2(r−1):指向每一层的索引块的指针占用\\
        2(r−1)B:每一层的索引块占用\\
        B-1:冗余未使用的数据块大小\\
        \item $O(r)$应该是最坏访问耗时，虽然是二重指针，但是确定第一个指针（去哪一层）可能需要$O(r)$\\
        例如要访问最后一个元素，你需要先整除$B^{r-1}$，发现大于该层所存数据块数，进入下一层，继续取模后整除$B^{r-2}$。。。直到整除$B$，一共r次
        \item 依旧没看懂信用什么的（TODO）
        \item 依旧没看懂第8部分，好像直接就用了神秘定理，但这个定理对不对，怎么证的是没说的（TODO）
    \end{enumerate}
\end{quote}

\noindent 2026年7月10日
\begin{quote}
    第一次通读原文

    第三次通读仍有一些问题，例如：
    \begin{enumerate}
        \item 第1部分中说可以通过拼接两个仅支持单端增删的数组实现Deque，然而这明显是错的，又或者不是显然可构造的\\
        两个仅支持单端增删的数组拼接确实可以很好地支持双端增删，但是如果我像队列那样，一边增，一边删呢？当删的那一端的数组没有元素时就会有问题\\
        所以可能有其他构造方法？又或者我对Deque理解有误？待议
        \item 8.2部分没看太懂
        \item Proof of Theorem8.6.部分中 $Note\ that\ if\ N>k\ and\ N\leq \binom{n+k}{k}\ ,then\ n\geq 2.$没看懂，n=1时不也成立吗

    \end{enumerate}
    对第二次阅读问题的解决：
    \begin{enumerate}
        \item 破案了，原文写的就是$\frac{1}{\alpha+1}$，然而原文疑似写错了，具体推论如下:\\
        如果只考虑增容：冗余占比应为$\frac{\alpha}{\alpha+1}$\\
        如果同时考虑缩容：冗余占比应大致为$\frac{2*\alpha+\alpha^2 }{\left( \alpha+1 \right)^2}$\
        \item 我觉得是主索引块指针占用空间为1，存储数组长度占用空间为1，总和为2
        \item 第3.3部分无平方根变体AI翻译部分出现删减，按原文构造方法是合理的
        \item 最多支持2B块很明显是缓冲缩容增容，如果只支持B块，那B块全满了之后增容一次，缩容两次，增容两次，缩容两次，增容两次，缩容两次。。。开销极大
        \item 我只能说豆包废物，靠后的理论章节删一大堆害我看的云里雾里的
    \end{enumerate}
    对第一次阅读问题的解决：
    \begin{enumerate}
        \item N+(2r−1)+2(r−1)B+B的组成修正如下:\\
        N:有效数据占用\\
        2r−1:指向每一层的索引块的指针占用,以及每一部分长度数据占用\\
        2(r−1)B:每一层的索引块占用\\
        B:冗余未使用的数据块大小\\
        实际上原文有图，十分直观，AI翻译文档无图导致产生疑问
        \item 已在第二次通读AI文档解决
        \item 已通过原文搞明白了（AI怎么删了这么多）
        \item 已通过原文搞明白了，除了上文说的8.2部分（AI怎么删了这么多）
    \end{enumerate}
    看原文又有一种局部看得懂，整合起来又有点晕晕的奇妙感觉\\

    趣事：第4部分证明定理4.2的时候，倒二行把N写成n，害我看了老半天（bad）\\
    第6.1部分惊现神秘不等式$for\ 1\ \leq\ r\ \leq\ r\ -\ 1$\\
    依旧是第6.1部分，在定理6.1的上一段，出现语句n1 is decremented and n1 is set to B.按理来说应该是n1 is decremented and n0 is set to B.（又害我看了老半天）\\
    看到这段话能绷住的也是神人了(6.3最后一段)：\\
    Obtaining a constant access time independent of r is only of theoretical interest as in practice the\\
    value of r used is expected to be very small. The simple O(logr) binary search algorithm, or even\\
    the na¨ıve O(r) algorithm, for locating the value of k would probably work best.\\
    你是说我费劲巴拉地看懂的你写的一大堆理论推理只是写来玩的，实践中根本没必要！？( $\xi =(´ο`*))$ 唉，理论佬)\\
    Theorem8.6.证明 n与n+1情况等价的时候又写错了（怎么这么多错误啊啊啊啊啊啊）\\
    Claim8.12.的两个case在同一个地方写错了我蚌埠住了（疑似复制粘贴）\\
    Lemma8.15.的最后一个式子把j写成i了\\

\end{quote}
\noindent 2026年7月12日
\begin{quote}
    之前阅读的是arXiv上的预发行版本，于是我找到了SIAM上正式发行的版本进行阅读(选择性阅读)
    第一件事就是看看我发现的那些错误有没有被修正:
    \begin{enumerate}
        \item 第4部分证明定理4.2的时候，倒二行把N写成n，仍未修正
        \item 第6.1部分第一段最后一句的神秘不等式$for\ 1\ \leq\ r\ \leq\ r\ -\ 1$仍未修正（推测应为$for\ 1\ \leq\ i\ \leq\ r\ -\ 1$）
        \item 第6.1部分，在定理6.1的上一段，语句n1 is decremented and n1 is set to B.仍未修正。很明显，如果要给n1赋值，那前一步的自减没有意义
        \item Theorem8.6.证明 n与n+1情况等价的错误已修正（arXiv版本误将$\binom{n+k}{k}$写作$\binom{n+k}{k+1}$）
        \item Claim8.12.的两个case的错误均已修正
        \item Lemma8.15.的最后一个式子误把j写成i的错误已修正
    \end{enumerate}
    看来只改了一半呢
    接下来是对第三次阅读问题的解答：
    \begin{enumerate}
        \item 依然不懂如何构造
        \item SIAM版本对8.2部分有补充解释，大致是看懂了
        \item Proof of Theorem8.6.部分中 $Note\ that\ if\ N>k\ and\ N\leq \binom{n+k}{k}\ ,then\ n\geq 2.$经思考后，认为极可能为新的书写错误，理由如下:\\
        先给出相关上下文 Note that if $N>k$ and $N\leq \dbinom{n+k}{k}$, then $n\geq 2$. Claim 8.10 shows that for every $\mathbf{a}\in Q_{N,k}$ we have $C(\mathbf{a})=(N+1)n - \dbinom{n+k}{k+1}$.\\
        Claim 8.10. If Theorem 8.6(i) holds for every $(N',k')<(N,k)$, and if $\binom{n+k-1}{k} - 1\leq N<\binom{n+k}{k} - 1$, for some $n\geq 1$ ......\\
        容易看出，原文一开始给N限定的范围$N\leq \dbinom{n+k}{k}$，比Claim 8.10.要求的范围$\binom{n+k-1}{k} - 1\leq N<\binom{n+k}{k} - 1$更宽松，然而原文第二句就直接运用Claim 8.10的结论，这显然是不合理的\\
        故修正后原文应为$Note\ that\ if\ N>k\ and\ N<\binom{n+k}{k} - 1\ ,then\ n\geq 2.$，这样比较合理一点。\\
        不过奇怪的是，$n\geq 2$这个结论并没有作用，因为Claim 8.10.仅要求$n\geq 1$，后文所使用的所有Claim也仅要求$n\geq 1$
    \end{enumerate}
\end{quote}


本文Latex格式有AI辅助修正，美化，但内容不变。

\end{document}
